\section{Introduction}

Electoral competitiveness is a fundamental democratic value: competitive
elections hold representatives accountable, encourage voter participation,
and reduce the influence of incumbency advantages.
Partisan gerrymandering directly reduces competitiveness by design ---
the map-drawer's explicit goal is to create safe districts for their party
while concentrating opposition voters into a small number of supermajority
districts.

Algorithmic redistricting cannot replicate this strategy because it cannot
access partisan data.
But the question of interest for policy is stronger: do algorithmic plans
produce \emph{more} competitive districts than enacted gerrymanders,
or do they merely avoid making things worse?

This paper provides the first systematic answer using five study states,
two plan types (bisect algorithmic and enacted 2022), and two competitiveness
metrics: the competitive district fraction (margin $< 10$ pp) and the
efficiency gap~\citep{stephanopoulos2015partisan}.

\paragraph{Causal identification.}
We treat the comparison between bisect and enacted plans as a
counterfactual: had the enacted plan been replaced by the algorithmic plan,
what fraction of competitive districts would exist?
This is not a randomised experiment; we cannot control for all factors that
differ between the two plans beyond competitiveness.
However, since both plans use the same census geography and population data,
the comparison isolates the effect of the mapmaking process itself.
